\documentclass[11pt,a4paper]{article}

\usepackage[margin=2.5cm]{geometry}
\usepackage{amsmath,amssymb,amsthm}
\usepackage{algorithm}
\usepackage{algpseudocode}
\usepackage{booktabs}
\usepackage{hyperref}
\usepackage{xcolor}
\usepackage{enumitem}
\usepackage{microtype}

\hypersetup{
  colorlinks=true,
  linkcolor=blue!70!black,
  citecolor=green!50!black,
  urlcolor=blue!70!black,
}

\newtheorem{definition}{Definition}
\newtheorem{remark}{Remark}

\title{\textbf{Constrained Path Planning in a Dynamic 2-D Cost Landscape}\\[0.4em]
\large Mathematical Problem Formulation and Proposed Solution}
\author{Path-Solver --- Technical Reference}
\date{}

\begin{document}
\maketitle
\tableofcontents
\bigskip

% ============================================================
\section{Overview}
% ============================================================

This document describes a constrained kinodynamic path-planning problem in a
2-D environment whose traversal cost evolves stochastically over time.
The planner must find a set of diverse, ranked candidate trajectories from a
fixed start pose to a goal region, respecting vehicle kinematic limits,
obstacle clearance requirements, and a monotone forward-progress condition.
When new environmental observations arrive (``trigger events''), the planner
reruns from scratch on the updated cost field and presents a refreshed ranked
list to the operator.

The code base consists of six Python modules:

\begin{center}
\begin{tabular}{ll}
\toprule
Module & Responsibility \\
\midrule
\texttt{cost\_landscape.py} & RBF cost field with OU-process dynamics \\
\texttt{vehicle\_model.py}  & Vehicle geometry and kinematic limits \\
\texttt{event\_emulator.py} & Poisson-clock trigger system \\
\texttt{graph\_search.py}   & Diverse A* seed generation \\
\texttt{optimizer.py}       & Constrained NLP path optimizer \\
\texttt{main.py}            & Orchestrator and visualizer \\
\bottomrule
\end{tabular}
\end{center}

All numeric values referenced symbolically in the body of this document are
consolidated in the Parameters Summary (Section~\ref{sec:params}).

% ============================================================
\section{Environment Model --- Dynamic Cost Landscape}
% ============================================================

\subsection{Scalar Cost Field}

The environment is a bounded 2-D domain
$\mathcal{D} = [x_{\min}, x_{\max}] \times [y_{\min}, y_{\max}]$.
A scalar cost field $C : \mathcal{D} \to \mathbb{R}_{>0}$ is defined as a
superposition of anisotropic Gaussian (Radial Basis Function, RBF) sources plus
a constant background term:

\begin{equation}
  \label{eq:cost_field}
  C(\mathbf{p}) = c_{\mathrm{bg}} + \sum_{s \in \mathcal{S} \cup \mathcal{O}}
    A_s \exp\!\left(
      -\frac{1}{2}\left[
        \left(\frac{u_s(\mathbf{p})}{\sigma_x^s}\right)^{\!2}
        + \left(\frac{v_s(\mathbf{p})}{\sigma_y^s}\right)^{\!2}
      \right]
    \right),
\end{equation}

where $\mathbf{p} = (x, y)^\top$, $c_{\mathrm{bg}}$ is a constant background
cost, $\mathcal{S}$ is the set of regular (soft) sources, and $\mathcal{O}$ is
the set of hard obstacle sources.

For source $s$ with center $\mathbf{c}_s = (c_x^s, c_y^s)^\top$ and
orientation angle $\vartheta_s$, the local rotated coordinates are

\begin{align}
  u_s(\mathbf{p}) &= \phantom{-}\cos(\vartheta_s)\,(x - c_x^s) + \sin(\vartheta_s)\,(y - c_y^s), \\
  v_s(\mathbf{p}) &= -\sin(\vartheta_s)\,(x - c_x^s) + \cos(\vartheta_s)\,(y - c_y^s).
\end{align}

\subsubsection*{Source parameters}

Each source $s \in \mathcal{S}$ is characterized by the parameter tuple
$(c_x^s, c_y^s, \sigma_x^s, \sigma_y^s, A_s, \vartheta_s)$ together with an
integer \emph{age} counter $k_s \ge 0$.
Initial values are drawn uniformly:
$A_s \sim \mathcal{U}(A_{\min}, A_{\max})$,
$\sigma_{\mathrm{base}} \sim \mathcal{U}(\sigma_{\min}^{(0)}, \sigma_{\max}^{(0)})$;
with probability $p_{\mathrm{elong}}$ the source is \emph{elongated}
($\sigma_y^s = \sigma_{\mathrm{base}} \cdot r$ with
$r \sim \mathcal{U}(r_{\mathrm{el,lo}}, r_{\mathrm{el,hi}})$)
and otherwise roughly isotropic
($\sigma_y^s = \sigma_{\mathrm{base}} \cdot r$ with
$r \sim \mathcal{U}(r_{\mathrm{iso,lo}}, r_{\mathrm{iso,hi}})$).

\subsubsection*{Obstacle sources}

Hard obstacles are placed uniformly in the interior
$[x_{\min} + \delta_{\mathrm{obs}},\; x_{\max} - \delta_{\mathrm{obs}}]^2$
with fixed parameters $A_{\mathrm{obs}}$ and
$\sigma_x = \sigma_y = \sigma_{\mathrm{obs}}$, yielding a cost that vastly
exceeds any traversal budget wherever the vehicle footprint overlaps the
obstacle center.
There are $N_{\mathrm{obs}}$ obstacles at initialization.

\subsection{Ornstein--Uhlenbeck Dynamics}

At each discrete time step $t$, the cost field evolves through an
Ornstein--Uhlenbeck (OU) update applied independently to the parameters of
every regular source $s \in \mathcal{S}$.

\subsubsection*{Amplitude}

\begin{equation}
  \label{eq:ou_amplitude}
  A_s \;\leftarrow\; A_s
    + \underbrace{\theta_A \,\bigl(\bar{A} - A_s\bigr)}_{\text{mean reversion}}
    + \underbrace{\sigma_A \,\xi_A}_{\text{volatility}},
  \qquad
  \xi_A \sim \mathcal{N}(0,1),
\end{equation}

with mean-reversion rate $\theta_A$, volatility $\sigma_A$, and
long-run mean $\bar{A} = (A_{\min} + A_{\max})/2$.
The result is clamped to $[A_{\mathrm{clamp}}, \infty)$.

\subsubsection*{Position}

\begin{align}
  c_x^s &\leftarrow c_x^s
    + \alpha_P\,\theta_P\,\bigl(\bar{x}_{\mathcal{D}} - c_x^s\bigr)
    + \sigma_P\,\xi_x, \\
  c_y^s &\leftarrow c_y^s
    + \alpha_P\,\theta_P\,\bigl(\bar{y}_{\mathcal{D}} - c_y^s\bigr)
    + \sigma_P\,\xi_y,
\end{align}

with mean-reversion rate $\theta_P$, drift volatility $\sigma_P$,
weak-reversion scaling factor $\alpha_P$, domain-center means
$\bar{x}_{\mathcal{D}} = (x_{\min} + x_{\max})/2$,
$\bar{y}_{\mathcal{D}} = (y_{\min} + y_{\max})/2$,
and $\xi_x, \xi_y \sim \mathcal{N}(0,1)$.
Positions are clamped to $\mathcal{D}$.
The small factor $\alpha_P$ makes the mean reversion of position very weak,
so sources perform nearly free Brownian drift.

\subsubsection*{Width}

\begin{equation}
  \sigma_x^s \leftarrow \operatorname{clip}\!\bigl(
    \sigma_x^s + \sigma_W\,\xi_{\sigma x},\;\sigma_W^{\min},\;\sigma_W^{\max}\bigr),
  \qquad
  \sigma_y^s \leftarrow \operatorname{clip}\!\bigl(
    \sigma_y^s + \sigma_W\,\xi_{\sigma y},\;\sigma_W^{\min},\;\sigma_W^{\max}\bigr),
\end{equation}

with width volatility $\sigma_W$ and clamp bounds $[\sigma_W^{\min}, \sigma_W^{\max}]$.

\subsubsection*{Orientation}

\begin{equation}
  \vartheta_s \leftarrow \vartheta_s + \sigma_\theta\,\xi_\theta, \qquad
  \xi_\theta \sim \mathcal{N}(0,1),
\end{equation}

with orientation drift volatility $\sigma_\theta$.

\subsubsection*{Obstacle drift}

Hard obstacle centers perform unconstrained Brownian drift without any OU
mean reversion or death:

\begin{equation}
  c_{x/y}^{\mathrm{obs}} \leftarrow c_{x/y}^{\mathrm{obs}}
    + \sigma_{\mathrm{obs,drift}}\,\xi, \quad \xi \sim \mathcal{N}(0,1),
\end{equation}

clamped to $[x_{\min} + \delta_{\mathrm{bd}},\; x_{\max} - \delta_{\mathrm{bd}}]$
(and likewise in $y$).

\subsection{Birth--Death Process}

At each update step, after the OU evolution, the source population undergoes:

\textbf{Death.}  Source $s$ is removed with probability
\begin{equation}
  p_{\mathrm{death}}(s) = \lambda_D \,\bigl(1 + \alpha_D\,k_s\bigr),
\end{equation}
where $\lambda_D$ is the base death rate, $\alpha_D$ is the age-scaling factor,
and $k_s$ is the source age.
Death is suppressed if the surviving count would fall below $N_{\min}$.

\textbf{Birth.}  A Poisson-distributed number of new sources is created:
\begin{equation}
  N_{\mathrm{new}} \sim \operatorname{Poisson}(\lambda_B),
\end{equation}
capped so the total count remains $\le N_{\max}$.
Newborn sources are initialized with $A_s$ scaled by $\alpha_{\mathrm{birth}}$
(attenuated at birth) and $k_s = 0$.

% ============================================================
\section{Vehicle Model}
% ============================================================

\subsection{Geometry}

The vehicle is modeled as a rectangular rigid body with physical dimensions

\begin{center}
\begin{tabular}{lll}
\toprule
Parameter & Symbol & Value \\
\midrule
Body length & $L$ & (see Table~\ref{sec:params}) \\
Body width  & $W$ & (see Table~\ref{sec:params}) \\
Half-length & $\ell = L/2$ & \\
Half-width  & $w = W/2$   & \\
\bottomrule
\end{tabular}
\end{center}

\subsection{Kinematic Limits}

The vehicle obeys the following bounds, used as hard constraints in the
optimizer (values in Section~\ref{sec:params}):

\begin{center}
\begin{tabular}{ll}
\toprule
Constraint & Symbol \\
\midrule
Maximum speed          & $v_{\max}$ \\
Minimum speed          & $v_{\min}$ \\
Maximum long.\ accel.  & $a_{\max}$ \\
Maximum lat.\ accel.   & $a_{\mathrm{lat,max}}$ \\
Maximum curvature      & $\kappa_{\max}$ \\
\bottomrule
\end{tabular}
\end{center}

\subsection{Oriented Bounding Box Footprint}

For obstacle-clearance evaluation the vehicle footprint at pose
$(x, y, \theta)$ is approximated by five sample points: the center and the
four corners of the oriented bounding box (OBB).
In the vehicle's local frame the five unit template points are

\begin{equation}
  \mathcal{F}_{\mathrm{unit}} =
  \{(0,0),\;(1,1),\;(1,-1),\;(-1,1),\;(-1,-1)\},
\end{equation}

scaled to world extents as $(\ell \cdot \hat{u},\; w \cdot \hat{v})$ and then
rotated and translated to world coordinates:

\begin{equation}
  \mathbf{p}_k = \mathbf{R}(\theta)\,
    \begin{pmatrix} \ell\,\hat{u}_k \\ w\,\hat{v}_k \end{pmatrix}
  + \begin{pmatrix} x \\ y \end{pmatrix},
  \quad
  \mathbf{R}(\theta) = \begin{pmatrix}
    \cos\theta & -\sin\theta \\ \sin\theta & \cos\theta
  \end{pmatrix}.
\end{equation}

The \emph{mean footprint cost} $\bar{C}_{\mathrm{fp}}$ and
\emph{max footprint cost} $\hat{C}_{\mathrm{fp}}$ are

\begin{equation}
  \bar{C}_{\mathrm{fp}}(x,y,\theta) = \frac{1}{5}\sum_{k=1}^{5} C(\mathbf{p}_k),
  \qquad
  \hat{C}_{\mathrm{fp}}(x,y,\theta) = \max_{k} C(\mathbf{p}_k).
\end{equation}

% ============================================================
\section{Trigger Event Model}
% ============================================================

A Poisson-process clock generates discrete ``trigger events'' that cause the
planner to rerun.
Inter-arrival times are drawn from an exponential distribution:

\begin{equation}
  \Delta t_k \sim \operatorname{Exp}\!\left(\frac{1}{\mu_T}\right),
\end{equation}

where $\mu_T$ is the mean trigger interval.
Between two consecutive events, the cost landscape is updated once via the
OU step described in Sections~2.2--2.3.
Each event records a cost-delta norm

\begin{equation}
  \|\Delta C\|_2 = \left\|
    C^{(t+1)}\big|_{\mathrm{grid}} - C^{(t)}\big|_{\mathrm{grid}}
  \right\|_2,
\end{equation}

evaluated on an $R_e \times R_e$ uniform sample grid with values clipped to
$[0, C_{\mathrm{ev,clip}}]$ before differencing.

% ============================================================
\section{Path Planning Problem}
% ============================================================

\subsection{State Representation}

A path is represented as a discrete sequence of $N$ \emph{waypoints}:

\begin{equation}
  \pi = \bigl\{(x_i, y_i, \theta_i, v_i)\bigr\}_{i=0}^{N-1},
\end{equation}

where $(x_i, y_i) \in \mathcal{D}$ is the position, $\theta_i \in [-\pi,\pi]$
is the heading, and $v_i > 0$ is the speed at waypoint $i$.

\subsection{Decision Variables}

The optimizer operates on a flat real vector of length $4N$:

\begin{equation}
  \mathbf{z} = \bigl[
    x_0,\, y_0,\, \theta_0,\, v_0,\;\;
    x_1,\, y_1,\, \theta_1,\, v_1,\;\;
    \ldots,\;\;
    x_{N-1},\, y_{N-1},\, \theta_{N-1},\, v_{N-1}
  \bigr]^\top \in \mathbb{R}^{4N}.
\end{equation}

\subsection{Objective Function}

The objective $J(\mathbf{z})$ has four additive terms:

\begin{equation}
  \label{eq:objective}
  J(\mathbf{z}) =
    \underbrace{
      \sum_{i=0}^{N-2} \bar{C}_{\mathrm{fp}}(x_i, y_i, \theta_i)\,\Delta s_i
    }_{J_{\mathrm{cost}}}
    + w_H \underbrace{\sum_{i=0}^{N-2} e_{H,i}^2}_{J_{\mathrm{heading}}}
    + w_S \underbrace{\sum_{i=0}^{N-2} \Delta\theta_i^2}_{J_{\mathrm{smooth}}}
    + w_V \underbrace{\sum_{i=0}^{N-2} \Delta v_i^2}_{J_{\mathrm{speed}}},
\end{equation}

where

\begin{align}
  \Delta s_i &= \|(x_{i+1} - x_i,\; y_{i+1} - y_i)\|_2, \\[4pt]
  e_{H,i}    &= \operatorname{wrap}_\pi\!\bigl(
                  \theta_i - \operatorname{atan2}(y_{i+1}-y_i,\, x_{i+1}-x_i)
                \bigr), \\[4pt]
  \Delta\theta_i &= \operatorname{wrap}_\pi(\theta_{i+1} - \theta_i), \\[4pt]
  \Delta v_i &= v_{i+1} - v_i,
\end{align}

and $w_H$, $w_S$, $w_V$ are regularization weights.
The operator $\operatorname{wrap}_\pi(\phi) = (\phi + \pi \bmod 2\pi) - \pi$
maps any angle to $(-\pi, \pi]$.

Term $J_{\mathrm{cost}}$ measures the arc-length-weighted integrated footprint
cost along the path.
$J_{\mathrm{heading}}$ penalizes misalignment between the declared heading
$\theta_i$ and the chord direction, coupling the orientation variable to the
spatial trajectory.
$J_{\mathrm{smooth}}$ penalizes heading discontinuities.
$J_{\mathrm{speed}}$ penalizes abrupt speed changes.

\subsection{Constraints}

\subsubsection*{Equality: Initial state}

The first waypoint is pinned to the vehicle's current state:

\begin{equation}
  x_0 = x_{\mathrm{init}}, \quad
  y_0 = y_{\mathrm{init}}, \quad
  \operatorname{wrap}_\pi(\theta_0 - \theta_{\mathrm{init}}) = 0, \quad
  v_0 = v_{\mathrm{init}}.
\end{equation}

\subsubsection*{Inequality: Goal disk}

The terminal waypoint must lie within the goal disk of radius $r_G$ centered at
$\mathbf{p}_G = (x_G, y_G)$:

\begin{equation}
  (x_{N-1} - x_G)^2 + (y_{N-1} - y_G)^2 \le r_G^2.
\end{equation}

\subsubsection*{Inequality: Monotone forward progress}

Each successive waypoint must reduce the Euclidean distance to the goal by at
most a slack $\varepsilon$:

\begin{equation}
  \label{eq:progress}
  d_i - d_{i+1} \ge -\varepsilon, \quad i = 0, \ldots, N-2,
\end{equation}

where $d_i = \|(x_i - x_G, y_i - y_G)\|_2$ and

\begin{equation}
  \varepsilon = \frac{\alpha_\varepsilon \cdot
    \|\mathbf{p}_G - \mathbf{p}_{\mathrm{init}}\|_2}{N}
\end{equation}

(unless overridden by the caller).
This constraint prevents path loops and back-tracking while tolerating small
lateral excursions.

\subsubsection*{Inequality: Speed bounds}

\begin{equation}
  v_{\min} \le v_i \le v_{\max}, \quad i = 0, \ldots, N-1.
\end{equation}

\subsubsection*{Inequality: Longitudinal acceleration}

Using the arc-length-parameterized approximation
$a_{\mathrm{long},i} \approx \bar{v}_i \cdot \Delta v_i / \Delta s_i$
with $\bar{v}_i = \tfrac{1}{2}(v_i + v_{i+1})$:

\begin{equation}
  \left| \bar{v}_i \cdot \frac{\Delta v_i}{\Delta s_i} \right| \le a_{\max},
  \quad i = 0, \ldots, N-2.
\end{equation}

\subsubsection*{Inequality: Path curvature}

Discrete curvature is approximated as $\kappa_i \approx |\Delta\theta_i / \Delta s_i|$:

\begin{equation}
  \kappa_i = \frac{|\operatorname{wrap}_\pi(\theta_{i+1} - \theta_i)|}{\Delta s_i}
  \le \kappa_{\max},
  \quad i = 0, \ldots, N-2.
\end{equation}

\subsubsection*{Inequality: Lateral acceleration}

\begin{equation}
  a_{\mathrm{lat},i} = \bar{v}_i^2 \,\kappa_i \le a_{\mathrm{lat,max}},
  \quad i = 0, \ldots, N-2.
\end{equation}

\subsubsection*{Inequality: Obstacle clearance}

The maximum footprint cost at every waypoint must stay below the hard threshold
$\tau$:

\begin{equation}
  \hat{C}_{\mathrm{fp}}(x_i, y_i, \theta_i) \le \tau,
  \quad i = 0, \ldots, N-1.
\end{equation}

\subsubsection*{Inequality: Domain boundary}

All waypoints must keep the vehicle OBB inside $\mathcal{D}$ by a margin
$m = \max(\ell, w)$:

\begin{equation}
  x_{\min} + m \le x_i \le x_{\max} - m, \quad
  y_{\min} + m \le y_i \le y_{\max} - m,
  \quad i = 0, \ldots, N-1.
\end{equation}

\subsubsection*{Summary}

The complete problem is a nonlinear program (NLP):

\begin{equation}
  \label{eq:nlp}
  \min_{\mathbf{z} \in \mathbb{R}^{4N}} J(\mathbf{z})
  \quad \text{subject to} \quad
  \mathbf{h}(\mathbf{z}) = \mathbf{0},\quad
  \mathbf{g}(\mathbf{z}) \ge \mathbf{0},
\end{equation}

where $\mathbf{h}$ collects the four initial-state equalities and $\mathbf{g}$
collects the $N-1$ progress constraints, $N-1$ speed-upper, $N$ speed-lower,
$2(N-1)$ acceleration, $N-1$ curvature, $N-1$ lateral-acceleration, $N$
obstacle-clearance, and $4N$ domain-boundary inequalities.

% ============================================================
\section{Solution Method}
% ============================================================

\subsection{Diverse Seed Generation via Penalty-Augmented A*}

Before invoking the NLP solver a set of $K$ topologically diverse seed paths
is constructed to initialize different local searches.
Algorithm~\ref{alg:seeds} describes the procedure.

\begin{algorithm}[H]
\caption{Diverse seed paths via penalty-augmented A*}
\label{alg:seeds}
\begin{algorithmic}[1]
\Require Cost landscape $C$, start $\mathbf{p}_s$, goal $\mathbf{p}_g$, $K$, $R$
\State Rasterize $C$ to $R \times R$ grid $G_{\mathrm{base}}$
\State $G_{\mathrm{pen}} \leftarrow \mathbf{0}_{R \times R}$; \quad seeds $\leftarrow []$
\For{$k = 1, \ldots, K$}
  \State $G_{\mathrm{cur}} \leftarrow G_{\mathrm{base}} + G_{\mathrm{pen}}$
  \State Run 8-connected A* on $G_{\mathrm{cur}}$; edge cost $= \delta_{\mathrm{world}} \cdot \bar{G}$
  \If{path $\pi_k$ found}
    \State Append $\pi_k$ to seeds
    \For{each point $\mathbf{q} \in \pi_k$}
      \State $G_{\mathrm{pen}} \mathrel{+}=
               A_{\mathrm{pen}} \exp\!\bigl(-\|\mathbf{p}_{\mathrm{grid}} - \mathbf{q}\|^2 /
               (2\sigma_{\mathrm{pen}}^2)\bigr)$
    \EndFor
  \EndIf
\EndFor
\State \Return seeds
\end{algorithmic}
\end{algorithm}

The A* edge cost between adjacent cells $(r,c)$ and $(r',c')$ is

\begin{equation}
  w(e) = \delta_{\mathrm{world}} \cdot
         \max\!\left(\frac{G_{\mathrm{cur}}(r,c) + G_{\mathrm{cur}}(r',c')}{2},\;
         \varepsilon_{\mathrm{num}}\right),
\end{equation}

where $\delta_{\mathrm{world}}$ is the Euclidean world-frame distance between
the two cell centers ($1$ or $\sqrt{2}$ grid steps for cardinal or diagonal
moves, respectively), and $\varepsilon_{\mathrm{num}}$ is a small positive
numerical floor that prevents zero edge weights.
Cells with cost $> \alpha_{\mathrm{block}} \cdot A_{\mathrm{obs}}$ are treated
as impassable.

The Gaussian repulsion penalty uses amplitude $A_{\mathrm{pen}}$ and
width $\sigma_{\mathrm{pen}}$, pushing subsequent searches into corridors that
are spatially separated from already-found paths.

\subsection{Arc-Length Resampling}

Each raw A* path (with variable number of grid-cell waypoints) is resampled to
exactly $N$ uniformly spaced points in arc-length:

\begin{equation}
  s_i^* = \frac{i}{N-1} \cdot L_{\mathrm{path}}, \quad i = 0, \ldots, N-1,
\end{equation}

where $L_{\mathrm{path}} = \sum_j \|\mathbf{p}_{j+1} - \mathbf{p}_j\|_2$ is
the total path length.
Linear interpolation between consecutive original points yields
$\tilde{\mathbf{p}}_i$ at target parameter $s_i^*$.

\subsection{NLP Initialization}

From a resampled seed $\{\tilde{x}_i, \tilde{y}_i\}$ the initial decision
vector $\mathbf{z}_0$ is built as follows.
Headings are set to consecutive chord angles:
$\theta_i^{(0)} = \operatorname{atan2}(\tilde{y}_{i+1} - \tilde{y}_i,\;
                                       \tilde{x}_{i+1} - \tilde{x}_i)$
(with $\theta_{N-1}^{(0)} = \theta_{N-2}^{(0)}$).
Speeds are initialized uniformly at $\alpha_v\,v_{\max}$, with
$v_0^{(0)} = v_{\mathrm{init}}$.

\subsection{Sequential Quadratic Programming (SLSQP)}

The NLP~\eqref{eq:nlp} is solved by the Sequential Least-Squares Quadratic
Programming (SLSQP) method as implemented in
\texttt{scipy.optimize.minimize(method=`SLSQP')}.
Solver settings:

\begin{center}
\begin{tabular}{ll}
\toprule
Setting & Symbol \\
\midrule
Maximum iterations & $I_{\max}$ \\
Function tolerance & $f_{\mathrm{tol}}$ \\
Gradient           & finite differences (default) \\
\bottomrule
\end{tabular}
\end{center}

SLSQP builds a local quadratic model of the Lagrangian at each iteration and
solves a quadratic program to obtain the search direction, then performs a
line search that respects the nonlinear constraints.

\subsection{Feasibility Check}

After the solver returns, a hard post-hoc feasibility check is applied with
tolerance $\delta_{\mathrm{tol}}$:

\begin{enumerate}[label=(\alph*)]
  \item $d_{N-1} \le r_G + \delta_{\mathrm{tol}}$ \quad (goal disk)
  \item $v_{\min} - \delta_{\mathrm{tol}} \le v_i \le v_{\max} + \delta_{\mathrm{tol}}$ for all $i$
  \item $d_i + \varepsilon + \delta_{\mathrm{tol}} \ge d_{i+1}$ for all $i$ \quad (progress)
  \item $\kappa_i \le \kappa_{\max} + \delta_{\mathrm{tol}}$ for all $i$ \quad (curvature)
  \item $a_{\mathrm{lat},i} \le a_{\mathrm{lat,max}} + \delta_{\mathrm{tol}}$ for all $i$
  \item $|a_{\mathrm{long},i}| \le a_{\max} + \delta_{\mathrm{tol}}$ for all $i$
  \item $\hat{C}_{\mathrm{fp}}(x_i, y_i, \theta_i) \le \tau + \delta_{\mathrm{tol}}$ for all $i$
  \item Domain bounds with margin $m$ (as above)
\end{enumerate}

A plan that passes all checks is labeled \textbf{feasible}; otherwise it is
\textbf{infeasible} (e.g., the solver converged to a local minimum violating a
constraint, or the problem has no monotone-progress solution on the current
landscape).

\subsection{Candidate Ranking and Labeling}

All $K$ solved plans are ranked by the composite key
$(\neg\,\mathrm{feasible},\; J_{\mathrm{cost}})$ --- feasible plans come first,
and among plans of equal feasibility status they are ordered by ascending
integrated footprint cost.
Plans are then labeled \textbf{Plan A}, \textbf{Plan B}, \ldots in rank order.

\subsection{Plan Metrics}

For each plan the following scalar metrics are reported:

\begin{align}
  J_{\mathrm{integrated}} &= \sum_{i=0}^{N-2} \bar{C}_{\mathrm{fp}}(x_i, y_i, \theta_i)\,\Delta s_i, \\[4pt]
  J_{\mathrm{peak}}       &= \max_{i=0,\ldots,N-1} \hat{C}_{\mathrm{fp}}(x_i, y_i, \theta_i), \\[4pt]
  L_{\mathrm{total}}      &= \sum_{i=0}^{N-2} \Delta s_i, \\[4pt]
  T_{\mathrm{total}}      &= \sum_{i=0}^{N-2} \frac{\Delta s_i}{\bar{v}_i}, \\[4pt]
  d_{\mathrm{goal}}       &= \|(x_{N-1} - x_G,\; y_{N-1} - y_G)\|_2.
\end{align}

% ============================================================
\section{Full Planning Loop}
% ============================================================

Algorithm~\ref{alg:main} gives the complete online planning loop executed by
the orchestrator.

\begin{algorithm}[H]
\caption{Online constrained path-planning loop}
\label{alg:main}
\begin{algorithmic}[1]
\Require $n_{\mathrm{events}}$, seeds $s_L$ (landscape), $s_E$ (emulator)
\State Initialize landscape $C$ with seed $s_L$
\State Initialize Poisson emulator with seed $s_E$
\State Define planning problem:
       start $(\mathbf{p}_s, \theta_{\mathrm{init}}, v_{\mathrm{init}})$,
       goal $\mathbf{p}_G$, radius $r_G$, waypoints $N$
\State \Call{Plan}{``Event 00 --- Initial''} \Comment{Plan on unperturbed landscape}
\For{each trigger event $e_k$, $k = 1, \ldots, n_{\mathrm{events}}$}
  \State Advance landscape via OU step (Sections~2.2--2.3)
  \State \Call{Plan}{``Event $k$''}
\EndFor
\Statex
\Procedure{Plan}{label}
  \State Generate $K$ diverse seed paths (Algorithm~\ref{alg:seeds})
  \State Resample each seed to $N$ waypoints
  \For{each seed $k$}
    \State Solve NLP~\eqref{eq:nlp} via SLSQP, initialized at $\mathbf{z}_0^{(k)}$
    \State Post-hoc feasibility check
    \State Compute plan metrics
  \EndFor
  \State Sort and label: Plan A, B, \ldots
  \State Save PNG visualization (cost heatmap + comparison table + speed profiles)
  \State Return ranked candidate list
\EndProcedure
\end{algorithmic}
\end{algorithm}

\textbf{Infeasible events.}
When the OU evolution places large obstacles across all monotone-progress
corridors it is mathematically impossible to satisfy
constraint~\eqref{eq:progress} while reaching the goal.
In this case all $K$ plans are correctly labeled infeasible; this serves as an
informative signal to the operator that no trajectory meeting all constraints
currently exists, and the system may request replanning after the next trigger
event.

% ============================================================
\section{Parameters Summary}
\label{sec:params}
% ============================================================

All numeric defaults are listed below.  Symbols used in the body of the
document reference entries in this table.

\subsection*{Domain and planning horizon}

\begin{center}
\begin{tabular}{llll}
\toprule
Symbol & Description & Value & Unit \\
\midrule
$x_{\min}, x_{\max}$ & Domain $x$ extent    & $0,\;100$ & m \\
$y_{\min}, y_{\max}$ & Domain $y$ extent    & $0,\;100$ & m \\
$N$                  & Waypoints per path   & $20$      & --- \\
$K$                  & Candidate seeds      & $3$       & --- \\
$\mathbf{p}_s$       & Start position       & $(10,10)$ & m \\
$\mathbf{p}_G$       & Goal position        & $(90,90)$ & m \\
$r_G$                & Goal radius          & $5$       & m \\
$\theta_{\mathrm{init}}$ & Initial heading  & $\pi/4$   & rad \\
$v_{\mathrm{init}}$  & Initial speed        & $5$       & m/s \\
\bottomrule
\end{tabular}
\end{center}

\subsection*{Vehicle model}

\begin{center}
\begin{tabular}{llll}
\toprule
Symbol & Description & Value & Unit \\
\midrule
$L$                    & Body length          & $5.0$   & m \\
$W$                    & Body width           & $2.5$   & m \\
$v_{\min}$             & Minimum speed        & $0.5$   & m/s \\
$v_{\max}$             & Maximum speed        & $15.0$  & m/s \\
$a_{\max}$             & Long.\ accel.\ limit & $3.0$   & m/s$^2$ \\
$a_{\mathrm{lat,max}}$ & Lat.\ accel.\ limit  & $2.0$   & m/s$^2$ \\
$\kappa_{\max}$        & Curvature limit      & $0.15$  & m$^{-1}$ \\
$\tau$                 & Obstacle cost threshold & $100$ & --- \\
\bottomrule
\end{tabular}
\end{center}

\subsection*{Cost landscape --- static field}

\begin{center}
\begin{tabular}{llll}
\toprule
Symbol & Description & Value & Unit \\
\midrule
$c_{\mathrm{bg}}$       & Background cost               & $0.1$        & --- \\
$A_{\min}, A_{\max}$    & Source amplitude range        & $0.5,\;5.0$  & --- \\
$\sigma_{\min}^{(0)}, \sigma_{\max}^{(0)}$ & Source sigma range & $3,\;15$ & m \\
$p_{\mathrm{elong}}$    & Elongation probability        & $0.3$        & --- \\
$r_{\mathrm{el,lo}}, r_{\mathrm{el,hi}}$   & Elongated sigma ratio range & $0.2,\;0.5$ & --- \\
$r_{\mathrm{iso,lo}}, r_{\mathrm{iso,hi}}$ & Isotropic sigma ratio range & $0.7,\;1.3$ & --- \\
$N_{\mathrm{obs}}$      & Number of hard obstacles      & $4$          & --- \\
$A_{\mathrm{obs}}$      & Obstacle amplitude            & $10^4$       & --- \\
$\sigma_{\mathrm{obs}}$ & Obstacle sigma ($x$ and $y$)  & $3.0$        & m \\
$\delta_{\mathrm{obs}}$ & Obstacle placement margin     & $15$         & m \\
\bottomrule
\end{tabular}
\end{center}

\subsection*{Cost landscape --- OU dynamics}

\begin{center}
\begin{tabular}{llll}
\toprule
Symbol & Description & Value & Unit \\
\midrule
$\theta_A, \sigma_A$          & OU amplitude rate / volatility & $0.3,\;0.4$   & --- \\
$A_{\mathrm{clamp}}$          & Minimum amplitude clamp        & $0.1$         & --- \\
$\theta_P, \sigma_P$          & OU position rate / volatility  & $0.1,\;1.5$   & m \\
$\alpha_P$                    & Weak position-reversion factor & $0.01$        & --- \\
$\sigma_W$                    & Width OU volatility            & $0.3$         & m \\
$\sigma_W^{\min}, \sigma_W^{\max}$ & Width clamp bounds        & $1,\;25$      & m \\
$\sigma_\theta$               & Orientation drift volatility   & $0.05$        & rad \\
$\sigma_{\mathrm{obs,drift}}$ & Obstacle drift std             & $0.3$         & m \\
$\delta_{\mathrm{bd}}$        & Obstacle boundary clamp margin & $5$           & m \\
\bottomrule
\end{tabular}
\end{center}

\subsection*{Cost landscape --- birth/death process}

\begin{center}
\begin{tabular}{llll}
\toprule
Symbol & Description & Value & Unit \\
\midrule
$\lambda_D$               & Base death rate            & $0.15$  & --- \\
$\alpha_D$                & Age factor in death prob.  & $0.1$   & --- \\
$\lambda_B$               & Birth (Poisson) rate       & $0.3$   & --- \\
$\alpha_{\mathrm{birth}}$ & Newborn amplitude scale    & $0.3$   & --- \\
$N_{\min}$                & Minimum source count       & $8$     & --- \\
$N_{\max}$                & Maximum source count       & $25$    & --- \\
\bottomrule
\end{tabular}
\end{center}

\subsection*{Trigger event emulator}

\begin{center}
\begin{tabular}{llll}
\toprule
Symbol & Description & Value & Unit \\
\midrule
$\mu_T$               & Mean trigger interval         & $8$    & s \\
$R_e$                 & Cost-delta sample grid size   & $50$   & --- \\
$C_{\mathrm{ev,clip}}$& Cost-delta clip ceiling       & $50$   & --- \\
\bottomrule
\end{tabular}
\end{center}

\subsection*{Optimizer}

\begin{center}
\begin{tabular}{llll}
\toprule
Symbol & Description & Value & Unit \\
\midrule
$w_H, w_S, w_V$          & Objective weights (heading, smooth, speed) & $10,\;2,\;0.5$ & --- \\
$\alpha_\varepsilon$      & Progress $\varepsilon$ fraction            & $0.1$          & --- \\
$\alpha_v$                & Speed initialization fraction              & $0.6$          & --- \\
$I_{\max}$                & SLSQP max iterations                       & $100$          & --- \\
$f_{\mathrm{tol}}$        & SLSQP function tolerance                   & $10^{-6}$      & --- \\
$\delta_{\mathrm{tol}}$   & Feasibility check tolerance                & $10^{-3}$      & --- \\
\bottomrule
\end{tabular}
\end{center}

\subsection*{Seed generator}

\begin{center}
\begin{tabular}{llll}
\toprule
Symbol & Description & Value & Unit \\
\midrule
$R$                      & A* grid resolution             & $80$  & --- \\
$A_{\mathrm{pen}}$       & Diversity penalty amplitude    & $5.0$ & --- \\
$\sigma_{\mathrm{pen}}$  & Diversity penalty width        & $8.0$ & m \\
$\alpha_{\mathrm{block}}$& Obstacle blocking fraction     & $0.01$& --- \\
\bottomrule
\end{tabular}
\end{center}

\end{document}
